\documentclass[11pt,a4paper]{article}
\usepackage[T1]{fontenc}
\usepackage[utf8]{inputenc}
\usepackage{lmodern,amsmath,amssymb,microtype}
\usepackage[margin=23mm]{geometry}
\usepackage{xcolor,hyperref,fancyhdr}
\definecolor{ink}{HTML}{173F4A}
\definecolor{muted}{HTML}{53636B}
\hypersetup{colorlinks=true,linkcolor=ink,urlcolor=ink,pdftitle={PSH Exhaustions: Research Catalogue}}
\pagestyle{fancy}
\fancyhf{}
\fancyhead[L]{\small\textsc{vibemath / PSH Exhaustions}}
\fancyhead[R]{\small Research catalogue}
\fancyfoot[L]{\small\color{muted}Draft collection \textbullet\ 23 October 2026}
\fancyfoot[R]{\small\thepage}
\setlength{\headheight}{14pt}
\setlength{\parindent}{0pt}
\setlength{\parskip}{6pt}
\setlength{\emergencystretch}{2em}
\newcommand{\entry}[2]{\par\medskip{\large\bfseries\color{ink}#1}\par\smallskip #2\par}
\begin{document}
\thispagestyle{plain}
{\small\scshape Several complex variables\hfill 23 October 2026}

{\LARGE\bfseries\color{ink}PSH Exhaustions}\par
{\Large Threefold counterexamples and surface criteria}\par
One manuscript \textbullet\ Editable current version \textbullet\ Archived snapshots

\medskip\hrule\medskip
This collection examines the distinction between continuous and smooth plurisubharmonic (psh) exhaustions. It packages a revised manuscript from the user's \emph{vibemath} research work, together with its earlier version and supplied revision materials.

\textbf{Current paper.} \emph{Continuous and Smooth Plurisubharmonic Exhaustions: Threefold Counterexamples and Partial Results for Surfaces}, by Zhijian Chen and Bo Xiao, dated 23 October 2026 (28 pages). The original 27-page revision and its 47-page predecessor are archived; the active source and PDF remain editable.

\entry{A. A K\"ahler threefold with a regularity gap}{
Theorem 1.1 constructs a smooth quasi-projective K\"ahler threefold
\[
M=\operatorname{Tot}(L^{-1})
\]
admitting a locally Lipschitz psh exhaustion but no $C^2$ psh exhaustion. The compact base is a ruled surface from Koike's modified Grauert example. The construction combines a continuous semipositive metric with an obstruction to every $C^2$ exhaustion.
}

\entry{B. The same gap on a non-K\"ahler threefold}{
Theorem 1.2 constructs, over any point $x\in M$, a proper modification $\mu:N\to M$ that is biholomorphic away from the fiber over $x$. The smooth threefold $N$ is non-K\"ahler and admits a locally Lipschitz psh exhaustion but no $C^2$ one. Hironaka's opposite-order blowups are an attributed input.
}

\entry{C. A line-bundle radius criterion}{
For a holomorphic line bundle $A$ over a compact complex manifold, Corollary 3.4 gives, for $r=2$ and $r=\infty$,
\[
\begin{gathered}
\operatorname{Tot}(A^{-1})\text{ has a }C^r\text{ psh exhaustion}\\
\Longleftrightarrow\\
A\text{ has a }C^r\text{ semipositive Hermitian metric}.
\end{gathered}
\]
Circle averaging and a high-level radius function give the reverse direction to the usual squared-norm construction.
}

\entry{D. Partial affirmative results for surfaces}{
Theorem 5.4 treats complements of nonzero effective divisors in nonprojective compact surfaces. Theorem 5.5 treats smooth $\mathbf P^1$-bundles with nonempty reduced simple normal crossing boundary. In both settings, continuous and smooth psh exhaustibility is equivalent to the absence of a negative definite connected boundary component.
}

\newpage
{\Large\bfseries\color{ink}Scope, dependencies, and reading guide}

\entry{What the surface section establishes}{
Proposition 5.2 shows that for a line bundle $A$ over a compact Riemann surface, continuous exhaustibility of $\operatorname{Tot}(A^{-1})$, nonnegative degree of $A$, and smooth exhaustibility are equivalent.

Proposition 5.6 gives further projective sufficient cases; Proposition 5.8 gives smooth square-zero boundary criteria. Corollary 5.7 supplies necessary numerical conditions for projective pairs remaining after these reductions. Those conditions do not construct a surface counterexample.

Table 1 concerns compact \emph{minimal} surface compactifiers under its explicit hypotheses. A change of compactifier can change the open complement. The table is not a classification of arbitrary noncompact surfaces.
}

\entry{Proof map and revision history}{
\textbf{Sections 2--3.} A cusp correction near a negative curve reaches the locally Lipschitz endpoint. Global gluing produces a metric; the radius criterion transfers the known $C^2$ metric obstruction to an exhaustion obstruction.

\textbf{Section 4.} Compactly supported modifications preserve the relevant exhaustion properties. A localized Hironaka configuration supplies the non-K\"ahler example.

\textbf{Section 5.} Boundary intersection matrices, collar constructions, surface classification, abundance, and neighborhood results yield the stated affirmative cases.

In the original surface revision, Sections 2--4 and the two principal theorem statements were unchanged from the predecessor. The archived snapshots, source hashes, patch, and Chinese change report preserve that record.
}

\entry{Attribution and verification status}{
The ruled-surface obstruction is credited to earlier work. Supporting criteria are not uniformly claimed as new; priority checking remains qualified.

Some surface branches invoke Theorem 1.6 of the cited September 2026 Chen--Filip--Sun--Tosatti--Zhang preprint. The supplied report records incomplete access to V\^aj\^aitu's published full text. These inherited qualifications remain in the repository.

Original hashes protect the archived snapshots; the active paper is checked for structure, references, and compilation. External journal review and formal verification are not established. The active source supplies authorship; licensing remains unassigned.
}

\entry{Questions left unresolved by this manuscript}{
The general continuous-to-$C^2$ exhaustion problem in complex dimension two is not resolved. Neither minimal counterexample dimension nor the optimal regularity threshold is determined. In particular, existence of a $C^1$ exhaustion on the constructed threefolds is not decided here.
}

\entry{Repository guide}{
Start with \texttt{CONTENTS.md} for precise statements and locations. The editable paper is in \texttt{preprints/}; original snapshots are in \texttt{archive/}. See \texttt{reviews/} for scope notes and revision materials, \texttt{reasoning\_traces/} for the editorial summary, and \texttt{UPLOAD\_TO\_GITHUB.md} for publishing instructions.
}

\smallskip
The catalogue structure is inspired by \href{https://github.com/openai/math}{openai/math}. This is a separate collection. The label \emph{vibemath} identifies the user's research context; it is not an author attribution.
\end{document}
